
Trustworthiness evaluations of large language models commonly treat adversarial fragility and miscalibration as positively coupled failure modes, implicitly assuming that a model brittle under adversarial attack will also be overconfident and poorly calibrated. This study tests that assumption directly for the adversarial fragility–calibration pair. A capability-controlled measure of adversarial fragility, termed Residual Instability (RI), is constructed by residualizing AdvGLUE accuracy drop against a PCA-derived capability index across 30 LLMs spanning 9 model families. Contrary to the coupled failure cascade prediction, RI significantly anticorrelates with Expected Calibration Error on the arc_challenge reasoning benchmark (Spearman partial ρ = −0.535, 95\% CI = [−0.782, −0.101], p = 0.0034, n = 30): models with higher adversarial fragility, after capability control, exhibit lower calibration error. Three mechanistic interpretations of this finding are considered, with a calibration–robustness trade-off driven by RLHF and instruction tuning identified as the most consistent with available evidence. Scope is explicitly limited to the RI–ECE dimension; sub-hypotheses addressing hallucination, safety, and output variance were not executed due to the gate failure on the primary mechanism. Several material limitations are noted, including that 73\% of AdvGLUE scores were OLS-estimated from 11 anchor models.

